% ==============================================================================
% secoes/00_resumo.tex - RESUMO E PALAVRAS-CHAVE
% Tema: A Contribuição de Federico Faggin para a Revolução dos Microprocessadores
%       e seus Impactos Sociais
% ==============================================================================

\begin{abstract}

Este artigo analisa o papel de Federico Faggin no desenvolvimento do Intel 4004 (1971), o primeiro microprocessador comercial em pastilha única de silício, e seus desdobramentos tecnológicos e sociais. Examina-se como a implementação da tecnologia MOS com porta autoalinhada de silício (\textit{Silicon Gate Technology}, SGT) e as soluções de projeto de circuito e leiaute viabilizaram a integração de uma unidade central de processamento em um único circuito integrado. A metodologia emprega triangulação entre fontes primárias, modelagem analítica e registros orais dos projetistas. Os resultados evidenciam que a supressão das capacitâncias parasitas de sobreposição pelo SGT, aliada a contatos enterrados e cargas dinâmicas com \textit{bootstrapping}, proporcionou densidade e velocidade indispensáveis para a integração monolítica da CPU. Discutem-se os desdobramentos sociotécnicos dessa transição, que substituiu circuitos dedicados por processamento programável, fundou os sistemas embarcados e reestruturou a automação fabril.
\end{abstract}

\begin{IEEEkeywords}
Federico Faggin, Intel 4004, \textit{Silicon Gate Technology}, Microprocessadores, Sistemas VLSI, Impactos Sociais.
\end{IEEEkeywords}
